% !TEX root = ../Tesis.tex
\chapter{Sitio de estudio}
El municipio de Salamanca está ubicado en el Estado de Guanajuato,
México ($20^\circ\,34'\,15''$\,N y $101^\circ\,11'\,50''$\,O) a
$1721$\,msnm (Figura~\ref{fig:mapa_salamanca}). Se encuentra en la zona
centro del Estado y colinda al Norte con los municipios de Dolores
Hidalgo y Allende; al Este con los municipios de Santa Cruz de
Juventino Rosas, Villagrán y Cortazar; al Oeste con Irapuato y Pueblo
Nuevo y finalmente al Sur colinda con los municipios de Valle de
Santiago y Jaral del Progreso. Cuenta con una superficie de
$757$\,km\textsuperscript{2} y una población de $273\,417$ habitantes
(Instituto de Ecología del Estado, 2022; Instituto Nacional de
Estadística y Geografía, 2010).

% AGREGAR MAPA
%\begin{figure}[H]
    %\centering
    %\includegraphics[width=0.85\textwidth]{figuras/mapa_salamanca.png}
    %\caption{Localización del municipio de Salamanca, Guanajuato,
    %México.}
    %\label{fig:mapa_salamanca}
%\end{figure}

La ciudad tiene un clima semicálido y ligeramente húmedo, con lluvias
de junio a agosto, cuyo rango de precipitación es de $600$ a
$900$\,mm. La temperatura anual promedio es de $18$\,°C, con altas
temperaturas en los meses de abril y junio, mientras que las
temperaturas más bajas son de diciembre a febrero. El suelo está
mayoritariamente enfocado a la agricultura, abarcando el $65.72\%$ del
territorio, mientras que la zona urbana es de $5.74\%$
(Campos-Ramos et al., 2012).

De acuerdo con un estudio realizado sobre la composición isotópica de
las precipitaciones de la Cuenca Alta del Río Laja, Guanajuato, en
2022, la línea meteórica local se representa por la
Ecuación~\ref{eq:lml} (Ramírez-González et al., 2024).

\begin{equation}
\delta\mathrm{D} = 7.349\,\delta^{18}\mathrm{O} + 0.693
\label{eq:lml}
\end{equation}

En el municipio de Salamanca se realiza la quema de esquilmos y se
operan ladrilleras. En estas últimas se usan diversos materiales como
la quema de basura, llantas, plásticos, diésel y madera. Como resultado
de esto, se generan partículas y gases que afectan la calidad del aire,
causando problemas a la salud pública. Los contaminantes más comunes
son óxidos de azufre (SO\textsubscript{x}), óxidos de nitrógeno
(NO\textsubscript{x}) y los compuestos orgánicos volátiles (COVs).
Estos compuestos afectan los ojos, irritan las fosas nasales y dañan la
piel, además de afectar el ecosistema donde se produce esta actividad
(Gobierno del Estado de Guanajuato y Secretaría de Seguridad Pública,
2023).

En el año 2021, el Estado de Guanajuato actualizó el inventario de
emisiones de contaminantes criterio y precursores. En él se reporta la
estimación de las emisiones de las cuatro principales fuentes de
emisión: fijas, de área, móviles y naturales. Los contaminantes PM10,
PM2.5, SO\textsubscript{2}, CO, NO\textsubscript{x} y ozono
troposférico son contaminantes criterio (Tabla~\ref{tab:contaminantes}).
Debido a que el ozono troposférico se forma a partir de reacciones
químicas que involucran óxidos de nitrógeno y compuestos orgánicos
volátiles (COVs) en presencia de luz solar, estos últimos se clasifican
como precursores de ozono (Secretaría de Agua y Medioambiente, 2021)
(Figura~\ref{fig:emisiones} a Figura~\ref{fig:emisiones_fin}).

\begin{table}[h!]
\centering
\small
\caption{Inventario de emisiones de contaminantes criterio y precursores
en el Estado de Guanajuato, 2021 (toneladas/año).}
\label{tab:contaminantes}
\resizebox{\textwidth}{!}{%
\begin{tabular}{lrrrrrrr}
\toprule
\textbf{Emisión} & \textbf{PM10} & \textbf{PM2.5} &
\textbf{SO\textsubscript{2}} & \textbf{CO} &
\textbf{NO\textsubscript{x}} & \textbf{COV} & \textbf{NH\textsubscript{3}} \\
\midrule
Emisiones totales
    & 2\,545.43 & 1\,799.66 & 7\,426.74 & 26\,974.73
    & 8\,650.06 & 14\,499.63 & 1\,178.08 \\
Fuentes fijas
    & 1\,478.50 & 1\,252.37 & 7\,331.79 & 2\,358.17
    & 2\,006.91 & 3\,016.05 & 107.75 \\
    & (36.23\%) & (36.03\%) & (84.76\%) & (37.51\%)
    & (29.07\%) & (49.11\%) & (60.58\%) \\
Fuentes de área
    & 772.91 & 280.23 & 43.05 & 1\,054.23
    & 81.90 & 3\,432.45 & 1\,046.92 \\
    & (4.00\%) & (3.18\%) & (12.29\%) & (2.15\%)
    & (3.47\%) & (4.01\%) & (2.96\%) \\
Fuentes móviles
    & 245.89 & 219.16 & 7.69 & 23\,411.99
    & 3\,820.73 & 2\,605.56 & 20.35 \\
    & (6.37\%) & (6.36\%) & (6.45\%) & (6.65\%)
    & (6.29\%) & (5.62\%) & (6.68\%) \\
Fuentes naturales
    & --- & --- & --- & ---
    & 2\,018.55 & 5\,441.65 & --- \\
    & & & & & (3.60\%) & (3.58\%) & \\
\bottomrule
\end{tabular}%
}
\end{table}

\subsection{Calidad del aire en Salamanca, Guanajuato}
A causa de la contingencia atmosférica, el gobierno de Guanajuato
implementó el programa ProAire en Salamanca, Celaya, Irapuato, León y
Silao, a lo que se le denominó ``Corredor Industrial''. Se basa en una
serie de acciones concretas para la reducción y control de las
diferentes emisiones provenientes de las fuentes con mayor
contribución a la contaminación.

Entre las medidas implementadas se encuentran la reforestación, la
construcción de ciclovías, el control de las cédulas de operación
anuales y el seguimiento a las industrias para reducir sus emisiones,
entre otras. De acuerdo con el avance de cumplimiento de las 37 medidas
establecidas en el ProAire, 32 se encuentran con un avance óptimo, 4
con retraso y una en peligro de ejecutarse, lo que corresponde a un
avance óptimo al 2019 del 86\% (Secretaría de Medio Ambiente y
Ordenamiento Territorial, 2022).

Asimismo, es posible sostener un control de las emisiones que emiten
las industrias instaladas en la zona. Estos datos se reportan a través
del Registro de Emisiones y Transferencia de Contaminantes (RETC). De
acuerdo con su giro industrial y pertinencia, pueden ser de
jurisdicción estatal (Tabla~\ref{tab:retc_estatal}) o federal
(Tabla~\ref{tab:retc_federal}).

\begin{table}[h!]
\centering
\footnotesize
\setlength{\tabcolsep}{4pt}
\caption{Empresas de jurisdicción estatal registradas en el RETC en
Salamanca, Guanajuato. Emisiones en toneladas/año.}
\label{tab:retc_estatal}
\resizebox{\textwidth}{!}{%
\begin{tabular}{lllrrr}
\toprule
\textbf{Clave} & \textbf{Nombre de la empresa o razón social} &
\textbf{Sector} & \textbf{CO\textsubscript{2}} &
\textbf{CH\textsubscript{4}} & \textbf{N\textsubscript{2}O} \\
\midrule
AEB325999112700 & ACEITES Y ESPECIALIDADES DEL BAJÍO, S.A. DE C.V. &
Industria química & 374\,607.98 & 5.99 & 0.61 \\
CMA339999112700 & CEBADAS Y MALTAS, S. DE R.L. DE C.V. &
Otras industrias manufactureras & 141.34 & 0.01 & 0 \\
DTM327111112700 & DAL-TILE MÉXICO, S. DE R.L. DE C.V. &
Fabricación de productos a base de minerales no metálicos
& 100\,012\,049.4 & 2\,488.50 & 309.28 \\
FEP324120112700 & FABRICACIÓN DE EMULSIONES Y PAVIMENTOS DE GUERRERO,
S.A. DE C.V. &
Fabricación de productos derivados del petróleo y del carbón
& 210\,306.18 & 3.34 & 0.33 \\
ACE336320112700 & FUJIKURA AUTOMOTIVE MÉXICO QUERÉTARO, S.A. DE C.V. &
Fabricación de equipo de transporte & 36\,036.00 & 0.92 & 0.15 \\
HPI311212112700 & HARINERA LOS PIRINEOS, S.A. DE C.V. &
Industria alimentaria & 27\,127.95 & 0.45 & 0.05 \\
ITI327122112700 & IBÉRICA TILES, S.A.P.I. DE C.V. &
Fabricación de productos a base de minerales no metálicos
& 74\,133\,481.99 & 1\,324.25 & 133.12 \\
JCS336390112700 & J CLIMA SISTEMAS MÉXICO, S.A. DE C.V. &
Fabricación de equipo de transporte & 10\,180.74 & 0.17 & 0.02 \\
PME336390112700 & PENSTONE MEXICANA, S.A. DE C.V. &
Fabricación de equipo de transporte & 162\,872.34 & 2.58 & 0.26 \\
PFR311411112000 & PRODUCTOS FRUGO, S.A. DE C.V. &
Industria alimentaria & 1\,990\,325.54 & 15.24 & 14.72 \\
PPV326110112700 & PRODUCTOS PLÁSTICOS VAZGÓN, S.A. DE C.V. &
Industria del plástico y del hule & --- & --- & --- \\
RSM339999112700 & SATUS AGER MÉXICO, S. DE R.L. DE C.V. &
Otras industrias manufactureras & 406\,295.30 & 7.21 & 0.72 \\
SSM331220112700 & SERVILAMINA SUMMIT MEXICANA, S.A. DE C.V. &
Industrias metálicas básicas & 39\,872.31 & 0.63 & 0.06 \\
TSA327999112700 & TRITURADOS SALAMANCA, S.A. DE C.V. &
Fabricación de productos a base de minerales no metálicos
& 115\,465.11 & 4.67 & 0.93 \\
\bottomrule
\end{tabular}%
}
\end{table}

\begin{table}[h!]
\centering
\footnotesize
\setlength{\tabcolsep}{4pt}
\caption{Empresas de jurisdicción federal registradas en el RETC en
Salamanca, Guanajuato. Emisiones en kg/año.}
\label{tab:retc_federal}
\resizebox{\textwidth}{!}{%
\begin{tabular}{llllrr}
\toprule
\textbf{Clave} & \textbf{Nombre de la empresa o razón social} &
\textbf{Sector} & \textbf{Sustancia} &
\textbf{Aire} & \textbf{Agua} \\
\midrule
PLO1102700112 & PEMEX Logística --- Terminal de Almacenamiento y
Despacho Salamanca & Petróleo y petroquímica &
CO\textsubscript{2} & 316\,965.1 & --- \\
PLO1102700117 & PEMEX Logística & Petróleo y petroquímica &
CO\textsubscript{2} & 1\,194\,300 & --- \\
PTI1102700059 & PEMEX Transformación Industrial, Refinería Salamanca
(Ing. Antonio M. Amor) & Petróleo y petroquímica &
Arsénico (polvos, respirables, vapores o humos) & 338.522 & --- \\
PTI1102700059 & PEMEX Transformación Industrial, Refinería Salamanca &
Petróleo y petroquímica & Formaldehído & 9\,313.302 & --- \\
PTI1102700059 & PEMEX Transformación Industrial, Refinería Salamanca &
Petróleo y petroquímica & Fenol & --- & 558.955 \\
PTI1102700059 & PEMEX Transformación Industrial, Refinería Salamanca &
Petróleo y petroquímica &
Cadmio (polvos, respirables, vapores o humos) & 102.069 & --- \\
PTI1102700059 & PEMEX Transformación Industrial, Refinería Salamanca &
Petróleo y petroquímica & Plomo (compuestos solubles) & 387.248 & --- \\
PTI1102700059 & PEMEX Transformación Industrial, Refinería Salamanca &
Petróleo y petroquímica & Cromo (compuestos solubles) & 216.705 & --- \\
PTI1102700059 & PEMEX Transformación Industrial, Refinería Salamanca &
Petróleo y petroquímica & Acetaldehído & 16\,678.969 & --- \\
TPL1102700153 & Tractocombustibles Plaza, S.A. de C.V. &
Petróleo y petroquímica & Tolueno & 4\,232.01 & --- \\
TPL1102700153 & Tractocombustibles Plaza, S.A. de C.V. &
Petróleo y petroquímica & Xileno (mezcla de isómeros) & 4\,584.93 & --- \\
TPL1102700153 & Tractocombustibles Plaza, S.A. de C.V. &
Petróleo y petroquímica & Benceno & 753.91 & --- \\
CGI1102700159 & CFE Generación I, Central Termoeléctrica Salamanca &
Generación de energía eléctrica & Bifenilos policlorados & --- & 4.073 \\
CGI1102700202 & CFE Generación I, Central Ciclo Combinado Cogeneración
Salamanca & Generación de energía eléctrica & Metano & 131\,658.268 & --- \\
CGI1102700202 & CFE Generación I, Central Ciclo Combinado Cogeneración
Salamanca & Generación de energía eléctrica &
CO\textsubscript{2} & 1\,684\,026\,621 & --- \\
AEX1102700127 & Aralo Express, S.A. de C.V. & Otros &
CO\textsubscript{2} & 18\,900\,000 & --- \\
CBLLU1102711 & Carboquímica Block, S.A. de C.V. &
Tratamiento de residuos peligrosos & CO\textsubscript{2} &
109\,168.952 & --- \\
GTK1102700007 & Y-Tec Keylex México, S.A. de C.V. & Automotriz &
CO\textsubscript{2} & 19\,818\,411 & --- \\
HCA5X1102711 & Henkel Capital, S.A. de C.V. & Química &
CO\textsubscript{2} & 4\,153\,217.6 & --- \\
ITI1102700197 & Ibérica Tiles, S.A.P.I. de C.V. &
Metalúrgica (incluye la siderúrgica) &
Plomo (compuestos solubles) & --- & 1.228 \\
MMM1102700005 & Mazda Motor Manufacturing de México, S.A. de C.V. &
Automotriz & Plomo (compuestos solubles) & --- & 55.984 \\
MMM1102700005 & Mazda Motor Manufacturing de México, S.A. de C.V. &
Automotriz & Mercurio (compuestos solubles) & --- & 1.501 \\
MMM1102700005 & Mazda Motor Manufacturing de México, S.A. de C.V. &
Automotriz & CO\textsubscript{2} & 12\,784\,377.65 & --- \\
MMM1102700005 & Mazda Motor Manufacturing de México, S.A. de C.V. &
Automotriz & Cromo (compuestos solubles) & --- & 131.36 \\
MMM1102700005 & Mazda Motor Manufacturing de México, S.A. de C.V. &
Automotriz & Cadmio (compuestos solubles) & --- & 10.25 \\
MMM1102700005 & Mazda Motor Manufacturing de México, S.A. de C.V. &
Automotriz & Arsénico (compuestos solubles) & --- & 2.793 \\
\bottomrule
\end{tabular}%
}
\end{table}

\subsection{Normas de calidad del aire}
En México, la NOM-172-SEMARNAT-2023 tiene por objeto establecer los
lineamientos para la obtención y comunicación del índice de la calidad
del aire. Lo anterior pretende brindar información estandarizada, clara
y oportuna, advirtiendo sobre los probables daños a la salud y las
medidas para reducir los impactos de la exposición (Secretaría de Medio
Ambiente y Recursos Naturales, 2024).

A partir de las concentraciones promedio horarias reportadas por las
estaciones de monitoreo, deben considerarse las cifras significativas
indicadas en la Tabla~\ref{tab:cifras_significativas} para cada
contaminante.

\begin{table}[h!]
\centering
\caption{Cifras significativas por contaminante según la
NOM-172-SEMARNAT-2023.}
\label{tab:cifras_significativas}
\begin{tabular}{llc}
\toprule
\textbf{Contaminante} & \textbf{Unidad} &
\textbf{Cifras significativas} \\
\midrule
PM10 & \si{\micro\gram\per\cubic\meter} & 0 \\
PM2.5 & \si{\micro\gram\per\cubic\meter} & 0 \\
Ozono (O\textsubscript{3}) & \si{\micro\gram\per\cubic\meter} & 3 \\
Dióxido de nitrógeno (NO\textsubscript{2}) & ppm & 3 \\
Dióxido de azufre (SO\textsubscript{2}) & ppm & 3 \\
Monóxido de carbono (CO) & ppm & 2 \\
\bottomrule
\end{tabular}
\end{table}

Asimismo, la normatividad mexicana establece los límites permisibles
para contaminantes atmosféricos, con la finalidad de contar con
indicadores que clasifiquen la calidad del aire e informen a la
población sobre los riesgos a la salud y las medidas de protección
(Tabla~\ref{tab:limites_permisibles}).

\begin{table}[H]
\centering
\footnotesize
\setlength{\tabcolsep}{4pt}
\caption{Límites permisibles para partículas suspendidas según la
normatividad mexicana.}
\label{tab:limites_particulas}
\resizebox{\textwidth}{!}{%
\begin{tabular}{lllrrr}
\toprule
\textbf{Contaminante} & \textbf{Indicador} & \textbf{Criterio} &
\textbf{Límite 1} & \textbf{Límite 2} & \textbf{Límite 3} \\
\midrule
PM10 & 24 horas &
Percentil 99 (frecuencia tolerada 1\% de veces al año) &
\SI{70}{\micro\gram\per\cubic\meter} &
\SI{60}{\micro\gram\per\cubic\meter} &
\SI{50}{\micro\gram\per\cubic\meter} \\
     & Anual &
Promedio anual de los promedios de 24 horas &
\SI{36}{\micro\gram\per\cubic\meter} &
\SI{28}{\micro\gram\per\cubic\meter} &
\SI{20}{\micro\gram\per\cubic\meter} \\
\midrule
PM2.5 & 24 horas &
Percentil 99 (frecuencia tolerada 1\% de veces al año) &
\SI{41}{\micro\gram\per\cubic\meter} &
\SI{33}{\micro\gram\per\cubic\meter} &
\SI{25}{\micro\gram\per\cubic\meter} \\
      & Anual &
Promedio anual de los promedios de 24 horas &
\SI{10}{\micro\gram\per\cubic\meter} &
\SI{10}{\micro\gram\per\cubic\meter} &
\SI{10}{\micro\gram\per\cubic\meter} \\
\bottomrule
\end{tabular}%
}
\end{table}

\begin{table}[H]
\centering
\footnotesize
\setlength{\tabcolsep}{4pt}
\caption{Límites permisibles para gases criterio según la normatividad
mexicana.}
\label{tab:limites_gases}
\resizebox{\textwidth}{!}{%
\begin{tabular}{llll}
\toprule
\textbf{Contaminante} & \textbf{Indicador} & \textbf{Criterio} &
\textbf{Límite} \\
\midrule
Ozono (O\textsubscript{3}) & 1 hora &
Máximo de los máximos diarios de las concentraciones horarias &
\SI{0.090}{ppm} \\
 & Promedio móvil de 8 horas &
Máximo de los máximos diarios de las concentraciones de los promedios
móviles &
\SI{0.065}{ppm} (1) \\
\midrule
Dióxido de azufre (SO\textsubscript{2}) & 1 hora &
Promedio aritmético de 3 años consecutivos de los percentiles 99
anuales del máximo diario de los datos horarios &
\SI{0.075}{ppm} \\
 & 24 horas &
Máximo de tres años consecutivos de los promedios de 24 horas &
\SI{0.040}{ppm} \\
\midrule
Dióxido de nitrógeno (NO\textsubscript{2}) & 1 hora &
Máximo de las concentraciones horarias &
\SI{0.106}{ppm} \\
 & Anual &
Promedio aritmético anual de las concentraciones horarias &
\SI{0.021}{ppm} \\
\midrule
Monóxido de carbono (CO) & 1 hora &
Máximo de las concentraciones horarias &
\SI{26}{ppm} \\
 & Promedio móvil de 8 horas &
Máximo de los promedios móviles de 8 horas &
\SI{9}{ppm} \\
\bottomrule
\end{tabular}%
}
\end{table}

\noindent (1) Los valores de O\textsubscript{3} para el promedio móvil
de 8 horas son \SI{0.065}{ppm}, \SI{0.060}{ppm} y \SI{0.051}{ppm},
correspondientes a los tres niveles de protección establecidos en la
norma.

